\section{Transformer 架构}

\subsection{Transformer Block}

Transformer 的核心构件是\textbf{Transformer Block}，由以下组件按顺序组成：

\begin{figure}[H]
\centering
\includegraphics[width=0.95\textwidth]{slides-images/slide-046.jpg}
\caption{Transformer Block 的结构：Multi-Head Self Attention $\rightarrow$ 残差连接 $\rightarrow$ Layer Norm $\rightarrow$ MLP $\rightarrow$ 残差连接 $\rightarrow$ Layer Norm\protect\footnotemark}
\end{figure}
\footnotetext{Slides 第47页。}

\begin{importantbox}{Transformer Block 的四个核心组件}
\begin{enumerate}
    \item \textbf{Multi-Head Self Attention}：让所有向量互相交互——``集体讨论''
    \item \textbf{MLP / FFN}：对每个向量独立进行非线性变换——``独立思考''
    \item \textbf{残差连接}（Residual Connection）：围绕 Self Attention 和 MLP 各加一个跳跃连接——改善梯度流动（与 ResNet 相同的原理）
    \item \textbf{Layer Normalization}：在残差连接之后进行归一化——稳定训练过程
\end{enumerate}

Self Attention 和 MLP \textbf{协同工作}：
\begin{itemize}
    \item Self Attention 负责\textbf{向量间}的信息交流
    \item MLP 负责\textbf{向量内}的非线性处理
\end{itemize}
\end{importantbox}

\subsection{完整 Transformer 模型}

一个完整的 Transformer 就是将多个 Transformer Block \textbf{堆叠}在一起：

\begin{figure}[H]
\centering
\includegraphics[width=0.95\textwidth]{slides-images/slide-047.jpg}
\caption{完整 Transformer 模型：多个 Transformer Block 的堆叠。原始 Transformer（2017）约 12 个 Block、2 亿参数；现代模型可达数百个 Block、数万亿参数\protect\footnotemark}
\end{figure}
\footnotetext{Slides 第48页。}

\begin{knowledgebox}{Transformer 的规模演变}
\begin{itemize}
    \item \textbf{2017}（原始 Transformer）：约 12 个 Block，$\sim$200M 参数
    \item \textbf{2020}（GPT-3）：96 个 Block，175B 参数
    \item \textbf{2023--2025}（现代大模型）：数百个 Block，数万亿参数
    \item 架构本身变化极小——核心的 Self Attention + MLP + 残差 + LayerNorm 结构始终未变
    \item 性能提升主要来自\textbf{更大的规模}和\textbf{更多的训练数据}
\end{itemize}
\end{knowledgebox}

\subsection{Transformer 用于语言建模}

将 Transformer 用于自回归语言建模：

\begin{enumerate}
    \item 将文本序列分割为 token
    \item 每个 token 通过 Embedding 层映射为向量
    \item 添加位置编码
    \item 通过多层 Transformer Block（带因果掩码的 Self Attention）
    \item 最终层输出经过线性投影 + softmax，预测下一个 token 的概率分布
\end{enumerate}

\subsection{Vision Transformer (ViT)}

\begin{figure}[H]
\centering
\includegraphics[width=0.95\textwidth]{slides-images/slide-050.jpg}
\caption{Vision Transformer (ViT)：将图像切分为 patch，每个 patch 投影为向量，送入标准 Transformer 处理\protect\footnotemark}
\end{figure}
\footnotetext{Slides 第51页。}

Transformer 同样可以处理图像：

\begin{enumerate}
    \item 将图像划分为固定大小的\textbf{patch}（如 $16 \times 16$ 像素）
    \item 每个 patch 通过线性投影映射为向量
    \item 添加位置编码
    \item 通过标准 Transformer Block 处理
    \item 对输出向量进行池化，通过线性层预测类别
\end{enumerate}

\begin{importantbox}{一种架构，多种数据}
Transformer 的强大之处在于其\textbf{通用性}：
\begin{itemize}
    \item \textbf{语言}：token $\rightarrow$ embedding $\rightarrow$ Transformer
    \item \textbf{图像}：patch $\rightarrow$ linear projection $\rightarrow$ Transformer
    \item \textbf{音频}：频谱帧 $\rightarrow$ linear projection $\rightarrow$ Transformer
    \item \textbf{视频}：时空 patch $\rightarrow$ linear projection $\rightarrow$ Transformer
\end{itemize}
同一种架构，只需改变输入的``向量化''方式，就能应用于截然不同的数据模态。
\end{importantbox}

\subsection{本章小结}

Transformer 架构由 Self Attention + MLP + 残差连接 + Layer Norm 四个组件构成。它的设计简洁、高度可扩展，自 2017 年以来几乎统治了深度学习的所有领域。通过不同的输入编码方式，同一个 Transformer 架构可以处理语言、图像、音频等多种数据类型。

%% ========== 拓展阅读 ==========

\subsection{拓展阅读}

\begin{itemize}
    \item Vaswani et al.: \textit{Attention Is All You Need}, NeurIPS 2017 \url{https://arxiv.org/abs/1706.03762}
    \item Bahdanau et al.: \textit{Neural Machine Translation by Jointly Learning to Align and Translate}, ICLR 2015 \url{https://arxiv.org/abs/1409.0473}
    \item Dosovitskiy et al.: \textit{An Image is Worth 16x16 Words: Transformers for Image Recognition at Scale}, ICLR 2021 \url{https://arxiv.org/abs/2010.11929}
    \item Jay Alammar: \textit{The Illustrated Transformer} \url{https://jalammar.github.io/illustrated-transformer/}
    \item Lilian Weng: \textit{Attention? Attention!} \url{https://lilianweng.github.io/posts/2018-06-24-attention/}
\end{itemize}

%% ========== 总结与延伸 ==========

